\subsection{Extension of isomorphisms}

PROPOSITION 8. --- Let \(\Omega\) be an algebraically closed extension of a field K, E and F two subextensions of \(\Omega\) and \(u\) a K-isomorphism of E onto F. For a K-automorphism \(v\) of \(\Omega\) to exist which extends \(u\) , it is necessary and sufficient for \(\Omega\) to have the same transcendence degree over E and F.

The condition is clearly necessary.

Suppose then that \(\Omega\) has the same transcendence degree over E and F and let us choose a transcendence basis B of \(\Omega\) over E and a transcendence basis C of \(\Omega\) over F. Since B and C are equipotent, Prop. 2 (V, p.~106) shows that \(u\) extends to a K-isomorphism \(u'\) of E(B) onto F(C). Since \(\Omega\) is an algebraic closure of E(B) and F(C), the Cor. of V, p.~23 shows that \(u'\) extends to an automorphism \(v\) of \(\Omega\) .

COROLLARY 1. --- Let \(\Omega\) be an algebraically closed extension of a field K and E a subextension of \(\Omega\) . Every K-automorphism of E extends to a K-automorphism of \(\Omega\) .

COROLLARY 2. --- Let \(\Omega\) be an algebraically closed extension of a field \(K\) , \(E\) and \(F\) two subextensions of \(\Omega\) and \(u\) a K-isomorphism of \(E\) onto \(F\) . If the transcendence degree of \(E\) over \(K\) is finite (in particular if \(E\) is algebraic over \(K\) ) then there exists a K-automorphism of \(\Omega\) extending \(u\) .

Let us denote by n, \(d(\mathrm{E})\) and \(d(\mathrm{F})\) the transcendence degrees of E over K, of \(\Omega\) over E and of \(\Omega\) over F respectively. The existence of the K-isomorphism u shows that the transcendence degree of F over K is equal to n.~By the Cor. of Th. 4 the transcendence degree of \(\Omega\) over K is equal to \(d(\mathrm{E}) + n\) and also to \(d(\mathrm{F}) + n\) . Therefore (E, III, p.~28, Prop. 8) we have \(d(\mathrm{E}) = d(\mathrm{F})\) and we can apply Prop. 8.

PROPOSITION 9. --- Let K be a field and \(\Omega\) an algebraically closed extension of K. Suppose that \(\Omega\) is not algebraic over K; then the set of elements of \(\Omega\) which are transcendental over K is infinite. Moreover, if \(x\) and \(y\) are two elements of \(\Omega\) which are transcendental over K, then there exists an automorphism \(u\) of \(\Omega\) over K such that \(u(x) = y\) .

Since \(\Omega\) is not algebraic over \(\mathbf{K}\) , there exists an element \(x\) of \(\Omega\) transcendental over \(\mathbf{K}\) ; then the elements \(x^n\) ( \(n \in \mathbf{N}\) ) are distinct and transcendental over \(\mathbf{K}\) . Suppose that \(x\) and \(y\) are distinct and transcendental over \(\mathbf{K}\) ; by Prop. 2 (V, p.~106) there exists a K-isomorphism \(\tilde{u}\) of \(\mathbf{K}(x)\) onto \(\mathbf{K}(y)\) such that \(\tilde{u}(x) = y\) ; since \(\mathbf{K}(x)\) is of transcendence degree 1 over \(\mathbf{K}\) , Cor. 2 of Prop. 8 shows that \(\tilde{u}\) extends to a K-automorphism \(u\) of \(\Omega\) .

PROPOSITION 10. --- Let K be a field, Ω an algebraically closed extension of K and G the group of K-automorphisms of Ω. Let \(x \in \Omega\) .

\begin{enumerate}
\def\labelenumi{\alph{enumi})}
\item
  For \(x\) to be algebraic over \(K\) it is necessary and sufficient for the set of elements \(u(x)\) as \(u\) runs over \(G\) to be finite,
\item
  For \(x\) to be \(p\) -radical over \(\mathbf{K}\) it is necessary and sufficient that \(u(x) = x\) for all \(u \in \mathbf{G}\) .
\end{enumerate}

In particular if K is perfect, the set of invariants of the group G is equal to K. Suppose first that x is transcendental. By Prop. 9 the set T of elements of \(\Omega\) transcendental over K is infinite, and for each \(y \in T\) there exists \(u \in G\) such that \(u(x) = y\) . Hence the set of elements \(u(x)\) as u runs over G is infinite.

Suppose now that \(x\) is algebraic over \(\mathbf{K}\) and denote by \(f\) its minimal polynomial over \(\mathbf{K}\) ; the set of roots of \(f\) in \(\Omega\) is finite and for each \(u \in \mathbf{G}\) we have \(f(u(x)) = u(f(x)) = 0\) . Hence the set of elements \(u(x)\) as \(u\) runs over \(\mathbf{G}\) is finite. This proves \(a\) .

Let L be the set of elements y of \(\Omega\) such that \(u(y)=y\) for all \(u\in G\) , and let \(\bar{K}\) be the relative algebraic closure of K in \(\Omega\) . By what has been said, L is a subextension of \(\bar{K}\) over K. Further (Cor. 1 to Prop. 8) every K-automorphism of \(\bar{K}\) is the restriction to \(\bar{K}\) of an element of G. Assertion b) of Prop. 10 now follows from Cor. 3 of V, p.~53.

